\chapter{Représentation matricielle d'une application linéaire}
\label{app:representation-lineaire-matrice}

Soient $V$ et $W$ deux espaces vectoriels de dimensions finies sur un corps
$\mathbb{K}$, où $\mathbb{K}$ désigne ici $\mathbb{R}$ ou $\mathbb{C}$.
On note
\[
\mathcal{B} = (e_1,\dots,e_N)
\]
une base de $V$, et
\[
\mathcal{B}' = (e'_1,\dots,e'_M)
\]
une base de $W$.

Soit maintenant une application linéaire
\[
f : V \rightarrow W.
\]
Tout vecteur $v \in V$ se décompose de manière unique dans la base $\mathcal{B}$ :
\[
v = \sum_{i=1}^{N} \alpha_i e_i,
\]
où les coefficients $\alpha_i$ sont les coordonnées de $v$ dans la base $\mathcal{B}$.

Par linéarité de $f$, on a alors :
\[
f(v)
=
f\left(\sum_{i=1}^{N} \alpha_i e_i\right)
=
\sum_{i=1}^{N} \alpha_i f(e_i).
\]

Comme $f(e_i) \in W$, chaque vecteur $f(e_i)$ peut lui-même être décomposé dans la base
$\mathcal{B}'$. Il existe donc des scalaires $a_{j i} \in \mathbb{K}$ tels que :
\[
f(e_i)
=
\sum_{j=1}^{M} a_{j i} e'_j.
\]

En remplaçant cette expression dans la formule précédente, on obtient :
\[
f(v)
=
\sum_{i=1}^{N} \alpha_i
\left(
\sum_{j=1}^{M} a_{j i} e'_j
\right).
\]

En réorganisant les sommes selon les vecteurs de la base $\mathcal{B}'$, il vient :
\[
f(v)
=
\sum_{j=1}^{M}
\left(
\sum_{i=1}^{N} a_{j i}\alpha_i
\right)
e'_j.
\]

Si l'on note
\[
[f(v)]_{\mathcal{B}'}
=
\begin{pmatrix}
\beta_1 \\
\vdots \\
\beta_M
\end{pmatrix},
\qquad
[v]_{\mathcal{B}}
=
\begin{pmatrix}
\alpha_1 \\
\vdots \\
\alpha_N
\end{pmatrix},
\]
alors les coordonnées de $f(v)$ vérifient :
\[
\beta_j
=
\sum_{i=1}^{N} a_{j i}\alpha_i.
\]

On peut donc écrire sous forme matricielle :
\[
[f(v)]_{\mathcal{B}'}
=
A [v]_{\mathcal{B}},
\]
où
\[
A =
(a_{j i})_{\substack{1 \leq j \leq M \\ 1 \leq i \leq N}}
\in \mathbb{K}^{M \times N}.
\]

Autrement dit, la matrice $A$ représentant l'application linéaire $f$ est construite colonne
par colonne : sa $i$-ème colonne est le vecteur des coordonnées de $f(e_i)$ dans la base
$\mathcal{B}'$ :
\[
A
=
\begin{pmatrix}
| & | &        & | \\
[f(e_1)]_{\mathcal{B}'} & [f(e_2)]_{\mathcal{B}'} & \cdots & [f(e_N)]_{\mathcal{B}'} \\
| & | &        & |
\end{pmatrix}.
\]

Ainsi, toute application linéaire entre espaces vectoriels de dimensions finies peut être
représentée par une matrice, dès lors que l'on a choisi une base dans l'espace de départ
et une base dans l'espace d'arrivée.

Il est important de noter que cette matrice dépend du choix des bases. Une même application
linéaire peut donc avoir des représentations matricielles différentes selon les bases utilisées.